\documentclass[11pt,a4paper]{article}

% ── Codificación y fuentes ────────────────────────────────────────────────────
\usepackage[utf8]{inputenc}
\usepackage[T1]{fontenc}
\usepackage{lmodern}
\usepackage[spanish]{babel}

% ── Geometría y espaciado ─────────────────────────────────────────────────────
\usepackage[top=2.5cm, bottom=2.5cm, left=2.8cm, right=2.8cm]{geometry}
\usepackage{setspace}
\setstretch{1.15}
\usepackage{parskip}

% ── Tablas ────────────────────────────────────────────────────────────────────
\usepackage{booktabs}
\usepackage{tabularx}
\usepackage{array}
\usepackage{longtable}
\usepackage{enumitem}

% ── Colores y cajas ───────────────────────────────────────────────────────────
\usepackage[dvipsnames]{xcolor}
\usepackage{tcolorbox}
\tcbuselibrary{skins, breakable}

% ── Cabeceras y pies de página ────────────────────────────────────────────────
\usepackage{fancyhdr}
\usepackage{lastpage}
\pagestyle{fancy}
\fancyhf{}
\fancyhead[L]{\small\textbf{Informe de Evaluación} --- Física para Bioanalistas}
\fancyhead[R]{\small Ana Alcala}
\fancyfoot[C]{\small Página \thepage\ de \pageref{LastPage}}
\renewcommand{\headrulewidth}{0.4pt}

% ── Hiperenlaces ──────────────────────────────────────────────────────────────
\usepackage[colorlinks=true, linkcolor=NavyBlue, urlcolor=NavyBlue]{hyperref}

% ── Paleta de colores personalizados ─────────────────────────────────────────
\definecolor{desempeno_excelente}{HTML}{1A6B2F}
\definecolor{desempeno_bueno}{HTML}{2E5A9C}
\definecolor{desempeno_suficiente}{HTML}{8B6914}
\definecolor{desempeno_deficiente}{HTML}{C0392B}
\definecolor{desempeno_insuficiente}{HTML}{7B241C}
\definecolor{grisFondo}{HTML}{F5F5F5}
\definecolor{azulTitulo}{HTML}{1B2A6B}
\definecolor{verdeFortaleza}{HTML}{1A6B2F}
\definecolor{naranjaMejora}{HTML}{A04000}

% ── Estilos de cajas ─────────────────────────────────────────────────────────
\tcbset{
  cajaTitulo/.style={
    enhanced, breakable,
    colback=azulTitulo!8, colframe=azulTitulo,
    fonttitle=\bfseries\large, coltitle=white,
    attach boxed title to top left={yshift=-2mm, xshift=4mm},
    boxed title style={colback=azulTitulo},
    top=6pt, bottom=6pt, left=8pt, right=8pt,
  },
  cajaFortaleza/.style={
    enhanced, breakable,
    colback=verdeFortaleza!6, colframe=verdeFortaleza!60,
    fonttitle=\bfseries, coltitle=verdeFortaleza,
    top=4pt, bottom=4pt, left=8pt, right=8pt,
  },
  cajaMejora/.style={
    enhanced, breakable,
    colback=naranjaMejora!6, colframe=naranjaMejora!60,
    fonttitle=\bfseries, coltitle=naranjaMejora,
    top=4pt, bottom=4pt, left=8pt, right=8pt,
  },
  cajaRecomendacion/.style={
    enhanced, breakable,
    colback=grisFondo, colframe=gray!50,
    fonttitle=\bfseries, coltitle=black,
    top=4pt, bottom=4pt, left=8pt, right=8pt,
  },
}

% ─────────────────────────────────────────────────────────────────────────────
\begin{document}

% ══ PORTADA ══════════════════════════════════════════════════════════════════
\begin{center}
  {\Large\bfseries\color{azulTitulo} Universidad de Oriente/Núcleo Sucre --- Física para Bioanalistas}\\[4pt]
  {\large Informe de Evaluación Preliminar}\\[12pt]
  \rule{\linewidth}{1.2pt}\\[6pt]
  {\LARGE\bfseries Ana Alcala}\\[4pt]
  \rule{\linewidth}{0.6pt}
\end{center}

% ══ FICHA TÉCNICA ════════════════════════════════════════════════════════════
\vspace{4pt}
\begin{tcolorbox}[cajaTitulo, title=Datos del informe]
\begin{tabular}{@{}llll@{}}
  \textbf{Estudiante:}  & Ana Alcala  &
  \textbf{Fecha:}       & 2026-08-08 \\[4pt]
  \textbf{Archivo PDF:} & \texttt{Ana\_Alcala.pdf}  &
  \textbf{Páginas:}     & 5 \\
\end{tabular}
\end{tcolorbox}

% ══ RESULTADO GLOBAL ═════════════════════════════════════════════════════════
\vspace{6pt}
\begin{center}
  \begin{tcolorbox}[
    enhanced,
    colback=white, colframe=desempeno_excelente,
    width=0.62\linewidth,
    boxrule=2pt,
    halign=center, valign=center,
    top=8pt, bottom=8pt,
  ]
    \centering
    {\large\bfseries Puntaje sugerido}\\[6pt]
    {\Huge\bfseries\color{desempeno_excelente} 9.9/10}\\[6pt]
    {\large Nivel de desempeño: \textbf{\color{desempeno_excelente} Excelente}}
  \end{tcolorbox}
\end{center}

% ══ RESUMEN GENERAL ══════════════════════════════════════════════════════════
\section*{Resumen general}
El estudiante demuestra una comprensión excepcional de los principios físicos relacionados con líquidos y fenómenos de transporte, aplicándolos correctamente a escenarios de bioanálisis. Las explicaciones conceptuales son claras, precisas y bien fundamentadas. Los cálculos son exactos y el manejo de unidades es consistente, con una única omisión menor en la pregunta 4c.

% ══ FORTALEZAS Y ÁREAS DE MEJORA ════════════════════════════════════════════
\vspace{4pt}
\begin{minipage}[t]{0.48\linewidth}
  \begin{tcolorbox}[cajaFortaleza, title={Fortalezas}]
\begin{itemize}[leftmargin=1.5em, itemsep=2pt]
  \item Profunda comprensión conceptual de los fenómenos físicos (termorregulación, mecánica respiratoria, difusión, ósmosis, capilaridad).
  \item Habilidad para aplicar principios físicos a situaciones clínicas y biológicas, argumentando y justificando cada respuesta.
  \item Precisión en los cálculos numéricos y uso correcto de fórmulas.
  \item Claridad y coherencia en la argumentación y justificación de las respuestas.
  \item Excelente manejo de unidades en la mayoría de los problemas.
\end{itemize}
  \end{tcolorbox}
\end{minipage}
\hfill
\begin{minipage}[t]{0.48\linewidth}
  \begin{tcolorbox}[cajaMejora, title={Areas de mejora}]
\begin{itemize}[leftmargin=1.5em, itemsep=2pt]
  \item Atención a la inclusión de unidades en todas las respuestas finales de cálculos.
\end{itemize}
  \end{tcolorbox}
\end{minipage}

% ══ OBSERVACIONES POR PREGUNTA ═══════════════════════════════════════════════
\section*{Observaciones por pregunta}

\begin{longtable}{>{\bfseries}p{2.8cm} p{1.6cm} p{7.0cm} p{2.8cm}}
  \toprule
  \textbf{Pregunta} & \textbf{Puntaje} & \textbf{Evaluación} & \textbf{Errores detectados} \\
  \midrule
  \endhead
  \bottomrule
  \endfoot
    \textbf{Pregunta 1: Termorregulación y Eficiencia de la Sudoración} & 2.0/2.0 & La respuesta es excelente en todas sus partes. La justificación física de la mayor sudoración requerida en el paciente B es muy clara y precisa, enfocándose en el calor latente de vaporización y la ineficiencia de la sudoración que gotea. La explicación del efecto de la humedad relativa al 100\% es completa, abarcando tanto el mecanismo físico (saturación del aire, anulación del gradiente de presión de vapor) como la consecuencia clínica (riesgo de hipertermia). Los cálculos son correctos y las unidades están bien manejadas. & \textit{---} \\
    \midrule
    \textbf{Pregunta 2: Mecánica Respiratoria y Ley de Laplace} & 2.0/2.0 & La explicación de la Ley de Laplace y sus implicaciones en alvéolos de diferente tamaño sin surfactante es impecable, describiendo correctamente el flujo anómalo de aire y el colapso. La descripción del papel del surfactante es precisa, destacando su capacidad para variar la tensión superficial y equilibrar las presiones. El cálculo es correcto, con la fórmula, sustitución y resultado adecuados, incluyendo las unidades. & \textit{---} \\
    \midrule
    \textbf{Pregunta 3: Optimización en Hemodiálisis y Primera Ley de Fick} & 2.0/2.0 & La aplicación de la Primera Ley de Fick para analizar la eficiencia de difusión es correcta, con una clara explicación de la relación proporcional y el aumento porcentual de la eficiencia. La justificación del riesgo de reducir excesivamente el espesor de la membrana es pertinente y bien fundamentada en la resistencia mecánica y el riesgo de ruptura. El cálculo de la tasa de transporte es exacto, con todos los datos, fórmula y resultado correctos. & \textit{---} \\
    \midrule
    \textbf{Pregunta 4: Errores de Infusión Intravenosa y Ósmosis} & 1.9/2.0 & Las explicaciones de los efectos del agua destilada (hipotónica) y la solución de NaCl al 5\% (hipertónica) en los eritrocitos son conceptualmente correctas y detalladas, describiendo adecuadamente la hemólisis y la crenación. El cálculo de la presión osmótica es numéricamente correcto, con la fórmula y sustitución adecuadas, pero se omite la unidad final en la respuesta. & \textit{Omisión de unidades en el resultado final del cálculo de la presión osmótica (parte c).} \\
    \midrule
    \textbf{Pregunta 5: Capilaridad y Toma de Muestras Sanguíneas} & 2.0/2.0 & La explicación del comportamiento de los líquidos en tubos hidrofóbicos es correcta, relacionando las fuerzas de cohesión/adhesión con el ángulo de contacto y la depresión capilar. La justificación de la optimización con microcapilares de menor radio es clara, basándose en la Ley de Jurin (h $\propto$ 1/r). El cálculo de la altura de ascenso es preciso, con todos los pasos y unidades correctas. & \textit{---} \\
    \midrule
\end{longtable}

% ══ ERRORES DETECTADOS ═══════════════════════════════════════════════════════
\section*{Análisis de errores}

\subsection*{Errores conceptuales}
\textit{Ninguno identificado.}

\subsection*{Errores procedimentales}
\begin{itemize}[leftmargin=1.5em, itemsep=2pt]
  \item Omisión de unidades en el resultado final del cálculo de la presión osmótica (Pregunta 4c).
\end{itemize}

% ══ RECOMENDACIONES AL ESTUDIANTE ════════════════════════════════════════════
\section*{Retroalimentación para el estudiante}
\begin{tcolorbox}[cajaRecomendacion, title={Dirigido a: Ana Alcala}]
Tu desempeño en este examen es sobresaliente. Demuestras una comprensión profunda y una excelente capacidad para aplicar los principios de la física a contextos biológicos y clínicos. Tus explicaciones son claras, tus cálculos precisos y tu razonamiento lógico es muy sólido. La única área de mejora menor es asegurarte de incluir siempre las unidades en los resultados finales de tus cálculos, ya que esto es crucial en física y en aplicaciones clínicas. ¡Felicidades por tu excelente trabajo!
\end{tcolorbox}

% ══ NOTA PARA EL PROFESOR ════════════════════════════════════════════════════
\section*{Nota para el profesor}
\begin{tcolorbox}[
  enhanced, breakable,
  colback=yellow!8, colframe=yellow!60!black,
  fonttitle=\bfseries, coltitle=black,
  title={Observaciones para revisión manual},
  top=4pt, bottom=4pt, left=8pt, right=8pt,
]
El estudiante ha demostrado un dominio excepcional de la materia. Todas las preguntas conceptuales y de aplicación han sido respondidas con gran claridad y precisión, justificando adecuadamente los fenómenos físicos. Los cálculos son correctos en su totalidad. Se ha realizado una pequeña penalización por la omisión de las unidades en el resultado final de la Pregunta 4c, pero esto no desmerece el alto nivel de comprensión y aplicación mostrado. El examen refleja un desempeño de nivel 'Excelente'.
\end{tcolorbox}

% ══ PIE DE INFORME ════════════════════════════════════════════════════════════
\vspace{12pt}
\begin{center}
  \footnotesize\color{gray}
  Informe generado automáticamente por el sistema de evaluación con IA (gemini-2.5-flash) y soporte LaTex. \\
  Este documento es una evaluación \textbf{preliminar} para ser revisado por el profesor antes de ser entregado al 
  estudiante. \\
  Generado el 2026-08-08 \textbullet\ BIOANALISIS Sistema Académico de Evaluación.
  \end{center}

\end{document}
